\newpage
\phantomsection
\label{prf:q-dense-in-itself}
\LRAProofFor{cor:q-dense-in-itself}

\begin{remark*}[Return]
\hyperref[cor:q-dense-in-itself]{Return to Theorem}
\end{remark*}

\begin{theorem*}[$\mathbb{Q}$ Is Dense in Itself]
The subset $\mathbb{Q}$ is order dense in itself.
\end{theorem*}

\begin{proof}
\textbf{Professional Standard Proof.}~\\
To say that $\mathbb{Q}$ is order dense in itself means precisely that for
every $r,s\in\mathbb{Q}$ with $r<s$, there exists $t\in\mathbb{Q}$ such that
\[
r<t<s.
\]
But this is exactly the statement of
\hyperref[thm:between-any-two-rationals-is-a-rational]{Theorem~\ref*{thm:between-any-two-rationals-is-a-rational}}.
Therefore $\mathbb{Q}$ is dense in itself.
\end{proof}

\begin{remark*}[Proof structure]
This corollary is a direct repackaging of the preceding theorem.  The theorem
asserts that every rational interval $(r,s)$ contains another rational.  The
phrase “$\mathbb{Q}$ is dense in itself” is exactly the structural name for
that property.  So no new construction is needed; the proof is just an
unfolding of terminology.
\end{remark*}

\begin{dependencies}
\begin{itemize}
  \item \hyperref[thm:between-any-two-rationals-is-a-rational]{Theorem~\ref*{thm:between-any-two-rationals-is-a-rational}}
\end{itemize}
\end{dependencies}

\textbf{Detailed Learning Proof.}~\\
% MIGRATION TODO: Existing proof content preserved; detailed/professional companion body needs review.
TODO: Expand the proof into a detailed learning proof.

